% !PRINTED 652
\weberpage{652}
%<<CONT
is set, and by substituting the expressions (3) for $\alpha_i$ in (7):
\weq{9}{\lambda \omega_r = \overset{i}{\Sigma} \omega_i \overset{s}{\Sigma} a_{s,i} t_{s,r}.}

By the definition of the norm (§ 172), $N(\lambda)$ is therefore the
determinant formed from the coefficients
\[
\overset{s}{\Sigma} a_{s,i} t_{s,r},
\]
and by the multiplication theorem for determinants one can decompose
 it into
\weq{10}{N(\lambda) = A T}
where $A$ has the meaning (4), and $T$ is the determinant:
\weq{11}{T = \Sigma \pm t_{1,1} t_{2,2} \ldots t_{n,n}.}

Here $T$ is a homogeneous function of degree $n$ in the $t_i$,
whose coefficients are integral functions in $Z$. $A$ is itself
an integral function in $z$. One now proves, as in Bd. II, § 164,
that $T$ is a unit. For if this were not the case, then
$T$ would have some linear factor $z - c$, and by an elementary
theorem on determinants one could determine holomorphic integral
functionals $y_i$ in $z$, not all of them divisible by $z - c$, in
such a way that
\[
u_s = y_1 t_{s1} + \cdots + y_n t_{sn}
\]
becomes divisible by $z - c$ (one can even give the $y_i$ constant
% [sic] printed konstanten Koeffizienten, the accusative plural object of geben without article requires konstante Koeffizienten (slip in the German only; the English is unaffected)
coefficients, since they can be reduced to their residues with respect
 to $z - c$).

If one now sets
\[
\omega = \omega_1 y_1 + \omega_2 y_2 + \cdots + \omega_n y_n,
\]
then it follows from (7)
\[
\lambda \omega = \alpha_1 u_1 + \alpha_2 u_2 + \cdots + \alpha_n u_n,
\]
and from this it follows, since the $\alpha_i$ are divisible by $\lambda$ and the $u_i$ by $z - c$,
that $\omega$ is divisible by $z - c$, which contradicts the definition of the
minimal basis. Accordingly it follows from (10) that $A$
is the absolute norm of the functional $\lambda$:
\weq{12}{N_a(\lambda) = \Sigma \pm a_{11} a_{22} \ldots a_{nn}.}

For a prime functional $\pi$ we can easily find a basis.
Not all the elements $\omega_1, \omega_2, \ldots, \omega_n$ of a minimal basis
can be divisible by $\pi$, since by the linear form
\weq{13}{\omega = x_1 \omega_1 + x_2 \omega_2 + \cdots + x_n \omega_n}
%CONT>>
